\documentclass[12pt,a4paper]{article}
\usepackage[utf8]{inputenc}
\usepackage{amsmath,amssymb}
\usepackage[round]{natbib}
\usepackage{geometry}
\geometry{margin=2.5cm}

\title{\textbf{Supplementary Material~S1}\\[6pt]
Rankine critical-height derivation\\[6pt]
\large Accompanying: \textit{A multiplicative correction factor for USLE erodibility\\in Plinthosols with subsurface hydraulic impedance}}

\author{Luiz Diego Vidal Santos, Emersson Guedes da Silva, Marcos Oliveira de Santana,\\
Renisson Neponuceno de Ara\'{u}jo Filho, Henrique Aroeira Soares}

\date{}

\begin{document}
\maketitle

\section*{S1. Rankine critical-height derivation}

The critical vertical-cut height $H_c$ (Eq.~4 of the main text) is derived from Rankine earth-pressure theory under total-stress analysis with neutral pore pressure ($u = 0$), consistent with the determination of effective cohesion ($c$) and internal friction angle ($\phi$) by direct-shear testing on saturated samples (ASTM D3080). The formulation assumes a vertical bank geometry ($90^\circ$ inclination) and a planar failure surface.

For a vertical cut in a cohesive soil, the active earth-pressure coefficient is

\begin{equation}
K_a = \tan^2\!\left(45^\circ - \frac{\phi}{2}\right)
\end{equation}

\noindent and the critical height at which the tension zone extends to the full depth of the cut is

\begin{equation}
H_c = \frac{2c}{\gamma_{\mathrm{sat}}} \cdot \tan\!\left(45^\circ + \frac{\phi}{2}\right) + \frac{2c}{\gamma_{\mathrm{sat}} \cdot \tan\!\left(45^\circ + \frac{\phi}{2}\right)}
\end{equation}

\noindent where $c$ is the effective cohesion (kPa), $\phi$ the internal friction angle ($^\circ$) and $\gamma_{\mathrm{sat}}$ the saturated unit weight (kN/m$^3$). In the studied Plinthosol, $c = 13.02$~kPa, $\phi = 34.93^\circ$ and $\gamma_{\mathrm{sat}} = 19.93$~kN/m$^3$ yield $H_c = 3.35$~m.

The influence of positive pore pressures under field conditions is evaluated separately in the slope-stability analysis (Section~4.3 of the main text) through the parameter $r_u$. The adoption of Rankine theory for vertical cuts is a conservative simplification; for bank geometries with inclination below~$90^\circ$, limit-equilibrium methods such as Bishop's simplified method may yield less conservative factor-of-safety values \citep{fredlund_rahardjo_1993}, an aspect warranting verification in complementary work.

\bibliographystyle{sn-basic}
\bibliography{../referencias_artigos}

\end{document}
